\begin{thebibliography}{99}

\bibitem{busia2025tiny}
P. Busia, M. A. Scrugli, V. J.-B. Jung, L. Benini, and P. Meloni,
``A Tiny Transformer for Low-Power Arrhythmia Classification on
Microcontrollers,'' \emph{IEEE Transactions on Biomedical Circuits and
Systems}, vol. 19, no. 1, pp. 142--152, 2025,
doi: 10.1109/TBCAS.2024.3401858.

\bibitem{moody2001impact}
G. B. Moody and R. G. Mark,
``The Impact of the MIT-BIH Arrhythmia Database,''
\emph{IEEE Engineering in Medicine and Biology Magazine},
vol. 20, no. 3, pp. 45--50, 2001,
doi: 10.1109/51.932724.

\bibitem{goldberger2000physiobank}
A. L. Goldberger \emph{et al.},
``PhysioBank, PhysioToolkit, and PhysioNet: Components of a New Research
Resource for Complex Physiologic Signals,''
\emph{Circulation}, vol. 101, no. 23, pp. e215--e220, 2000,
doi: 10.1161/01.CIR.101.23.e215.

\bibitem{dechazal2004automatic}
P. de Chazal, M. O'Dwyer, and R. B. Reilly,
``Automatic Classification of Heartbeats Using ECG Morphology and Heartbeat
Interval Features,''
\emph{IEEE Transactions on Biomedical Engineering},
vol. 51, no. 7, pp. 1196--1206, 2004,
doi: 10.1109/TBME.2004.827359.

\bibitem{pellegrini2020latent}
L. Pellegrini, G. Graffieti, V. Lomonaco, and D. Maltoni,
``Latent Replay for Real-Time Continual Learning,''
in \emph{Proc. IEEE/RSJ International Conference on Intelligent Robots and
Systems}, 2020, pp. 10203--10209,
doi: 10.1109/IROS45743.2020.9341460.

\bibitem{ravaglia2021tinyml}
L. Ravaglia, M. Rusci, D. Nadalini, A. Capotondi, F. Conti, and L. Benini,
``A TinyML Platform for On-Device Continual Learning With Quantized Latent
Replays,''
\emph{IEEE Journal on Emerging and Selected Topics in Circuits and Systems},
vol. 11, no. 4, pp. 789--802, 2021,
doi: 10.1109/JETCAS.2021.3121554.

\bibitem{kirkpatrick2017overcoming}
J. Kirkpatrick \emph{et al.},
``Overcoming Catastrophic Forgetting in Neural Networks,''
\emph{Proceedings of the National Academy of Sciences},
vol. 114, no. 13, pp. 3521--3526, 2017,
doi: 10.1073/pnas.1611835114.

\bibitem{hu2022lora}
E. J. Hu, Y. Shen, P. Wallis, Z. Allen-Zhu, Y. Li, S. Wang, L. Wang,
and W. Chen,
``LoRA: Low-Rank Adaptation of Large Language Models,''
in \emph{International Conference on Learning Representations}, 2022.

\bibitem{douillard2022dytox}
A. Douillard, A. Ram\'e, G. Couairon, and M. Cord,
``DyTox: Transformers for Continual Learning With Dynamic Token Expansion,''
in \emph{Proc. IEEE/CVF Conference on Computer Vision and Pattern
Recognition}, 2022, pp. 9285--9295.

\bibitem{lin2022ondevice}
J. Lin, L. Zhu, W.-M. Chen, W.-C. Wang, C. Gan, and S. Han,
``On-Device Training Under 256KB Memory,''
in \emph{Advances in Neural Information Processing Systems},
vol. 35, 2022, pp. 22941--22954.

\bibitem{rebuffi2017icarl}
S.-A. Rebuffi, A. Kolesnikov, G. Sperl, and C. H. Lampert,
``iCaRL: Incremental Classifier and Representation Learning,''
in \emph{Proc. IEEE Conference on Computer Vision and Pattern Recognition},
2017, pp. 2001--2010.

\bibitem{tiwari2022gcr}
R. Tiwari, K. Killamsetty, R. Iyer, and P. Shenoy,
``GCR: Gradient Coreset Based Replay Buffer Selection for Continual
Learning,''
in \emph{Proc. IEEE/CVF Conference on Computer Vision and Pattern
Recognition}, 2022, pp. 99--108.

\bibitem{borsos2020coresets}
Z. Borsos, M. Mutn\'y, M. Tagliasacchi, and A. Krause,
``Coresets via Bilevel Optimization for Continual Learning and Streaming,''
in \emph{Advances in Neural Information Processing Systems},
vol. 33, 2020, pp. 14879--14890.


\bibitem{hu1997patient}
Y. H. Hu, S. Palreddy, and W. J. Tompkins,
``A Patient-Adaptable ECG Beat Classifier Using a Mixture of Experts
Approach,''
\emph{IEEE Transactions on Biomedical Engineering},
vol. 44, no. 9, pp. 891--900, 1997,
doi: 10.1109/10.623058.

\bibitem{kiranyaz2016realtime}
S. Kiranyaz, T. Ince, and M. Gabbouj,
``Real-Time Patient-Specific ECG Classification by 1-D Convolutional Neural
Networks,''
\emph{IEEE Transactions on Biomedical Engineering},
vol. 63, no. 3, pp. 664--675, 2016,
doi: 10.1109/TBME.2015.2468589.

\bibitem{llamedo2011heartbeat}
M. Llamedo and J. P. Mart\'inez,
``Heartbeat Classification Using Feature Selection Driven by Database
Generalization Criteria,''
\emph{IEEE Transactions on Biomedical Engineering},
vol. 58, no. 3, pp. 616--625, 2011,
doi: 10.1109/TBME.2010.2068048.

\bibitem{mar2011optimization}
T. Mar, S. Zaunseder, J. P. Mart\'inez, M. Llamedo, and R. Poll,
``Optimization of ECG Classification by Means of Feature Selection,''
\emph{IEEE Transactions on Biomedical Engineering},
vol. 58, no. 8, pp. 2168--2177, 2011,
doi: 10.1109/TBME.2011.2113395.

\bibitem{hannun2019cardiologist}
A. Y. Hannun, P. Rajpurkar, M. Haghpanahi, G. H. Tison, C. Bourn,
M. P. Turakhia, and A. Y. Ng,
``Cardiologist-Level Arrhythmia Detection and Classification in Ambulatory
Electrocardiograms Using a Deep Neural Network,''
\emph{Nature Medicine}, vol. 25, no. 1, pp. 65--69, 2019,
doi: 10.1038/s41591-018-0268-3.

\bibitem{greenwald1990svdb}
S. D. Greenwald,
``Improved Detection and Classification of Arrhythmias in Noise-Corrupted
Electrocardiograms Using Contextual Information,''
Ph.D. dissertation, Harvard--MIT Division of Health Sciences and Technology,
Massachusetts Institute of Technology, Cambridge, MA, USA, 1990.

\bibitem{ansi2012aami}
Association for the Advancement of Medical Instrumentation,
\emph{Testing and Reporting Performance Results of Cardiac Rhythm and
ST-Segment Measurement Algorithms}, ANSI/AAMI EC57:2012, 2012.

\bibitem{houlsby2019adapters}
N. Houlsby, A. Giurgiu, S. Jastrzebski, B. Morrone, Q. de Laroussilhe,
A. Gesmundo, M. Attariyan, and S. Gelly,
``Parameter-Efficient Transfer Learning for NLP,''
in \emph{Proc. International Conference on Machine Learning},
2019, pp. 2790--2799.

\bibitem{cai2020tinytl}
H. Cai, C. Gan, L. Zhu, and S. Han,
``TinyTL: Reduce Memory, Not Parameters for Efficient On-Device Learning,''
in \emph{Advances in Neural Information Processing Systems},
vol. 33, 2020, pp. 11285--11297.

\bibitem{wang2021tent}
D. Wang, E. Shelhamer, S. Liu, B. Olshausen, and T. Darrell,
``Tent: Fully Test-Time Adaptation by Entropy Minimization,''
in \emph{International Conference on Learning Representations}, 2021.

\bibitem{liang2020shot}
J. Liang, D. Hu, and J. Feng,
``Do We Really Need to Access the Source Data? Source Hypothesis Transfer for
Unsupervised Domain Adaptation,''
in \emph{Proc. International Conference on Machine Learning},
2020, pp. 6028--6039.

\bibitem{lin2017focal}
T.-Y. Lin, P. Goyal, R. Girshick, K. He, and P. Doll\'ar,
``Focal Loss for Dense Object Detection,''
in \emph{Proc. IEEE International Conference on Computer Vision},
2017, pp. 2980--2988.

\bibitem{cui2019classbalanced}
Y. Cui, M. Jia, T.-Y. Lin, Y. Song, and S. Belongie,
``Class-Balanced Loss Based on Effective Number of Samples,''
in \emph{Proc. IEEE/CVF Conference on Computer Vision and Pattern
Recognition}, 2019, pp. 9268--9277.

\bibitem{vaswani2017attention}
A. Vaswani, N. Shazeer, N. Parmar, J. Uszkoreit, L. Jones, A. N. Gomez,
{\L}. Kaiser, and I. Polosukhin,
``Attention Is All You Need,''
in \emph{Advances in Neural Information Processing Systems},
vol. 30, 2017, pp. 5998--6008.

\bibitem{li2018lwf}
Z. Li and D. Hoiem,
``Learning Without Forgetting,''
\emph{IEEE Transactions on Pattern Analysis and Machine Intelligence},
vol. 40, no. 12, pp. 2935--2947, 2018,
doi: 10.1109/TPAMI.2017.2773081.

\bibitem{chaudhry2019tiny}
A. Chaudhry, M. Rohrbach, M. Elhoseiny, T. Ajanthan, P. K. Dokania,
P. H. S. Torr, and M. Ranzato,
``On Tiny Episodic Memories in Continual Learning,''
\emph{arXiv:1902.10486}, 2019.

\bibitem{lopezpaz2017gem}
D. Lopez-Paz and M. Ranzato,
``Gradient Episodic Memory for Continual Learning,''
in \emph{Advances in Neural Information Processing Systems},
vol. 30, 2017, pp. 6467--6476.

\bibitem{zenke2017si}
F. Zenke, B. Poole, and S. Ganguli,
``Continual Learning Through Synaptic Intelligence,''
in \emph{Proc. International Conference on Machine Learning},
2017, pp. 3987--3995.

\bibitem{banbury2021mlperftiny}
C. Banbury \emph{et al.},
``MLPerf Tiny Benchmark,''
in \emph{Proc. Neural Information Processing Systems Datasets and Benchmarks
Track}, 2021.

\bibitem{david2021tflitemicro}
R. David \emph{et al.},
``TensorFlow Lite Micro: Embedded Machine Learning for TinyML Systems,''
in \emph{Proc. Machine Learning and Systems (MLSys)}, vol. 3, 2021,
pp. 800--811.


\bibitem{ruib2024tinyml}
M. R\"ub, P. Tuchel, A. Sikora, and D. Mueller-Gritschneder,
``A Continual and Incremental Learning Approach for TinyML On-Device Training
Using Dataset Distillation and Model Size Adaption,''
\emph{arXiv:2409.07114}, 2024.

\bibitem{replay2025survey}
H.-S. Park, H.-C. Chu, M.-K. Sung, C. Kim, J. Lee, D.-W. Kim, and J. Lee,
``Survey on Replay-Based Continual Learning and Empirical Validation on
Feasibility in Diverse Edge Devices Using a Representative Method,''
\emph{Mathematics}, vol. 13, no. 14, art. 2257, 2025,
doi: 10.3390/math13142257.

\bibitem{ecgfoundation2024survey}
Y. Han, V. Murino, X. Liu, X. Zhang, and C. Ding,
``A Systematic Review on Foundation Models for Electrocardiogram Analysis:
Initial Strides and Expansive Horizons,''
\emph{arXiv:2410.19877}, 2024.

\bibitem{zhang2026amrecl}
P. Zhang, J. Yang, Y. Yao, X. Ruan, Y. Chen, X. Wang, and Y. Chen,
``Adaptive Memory Replay for Plasticity-Aware Edge Continual Learning,''
\emph{Expert Systems with Applications}, in press, available online 2026;
assigned to vol. 331, art. 133393 (issue dated 15 December 2026),
doi: 10.1016/j.eswa.2026.133393.

\bibitem{ravaglia2020memlatacc}
L. Ravaglia, M. Rusci, A. Capotondi, F. Conti, L. Pellegrini, V. Lomonaco,
D. Maltoni, and L. Benini,
``Memory-Latency-Accuracy Trade-Offs for Continual Learning on a RISC-V
Extreme-Edge Node,'' in \emph{2020 IEEE Workshop on Signal Processing Systems
(SiPS)}, 2020, pp. 1--6, doi: 10.1109/SIPS50750.2020.9195220.

\bibitem{mitdb202note}
PhysioNet, ``MIT--BIH Arrhythmia Database, record 202 header (\texttt{202.hea}),''
MIT--BIH Arrhythmia Database v1.0.0, PhysioNet, 2005.
[Online]. Available: \url{https://physionet.org/content/mitdb/1.0.0/202.hea}

\bibitem{zheng2026ewclora}
Y. Zheng, Y. Zhang, J. van de Weijer, G. M. van de Ven, S. Du, X. Zhang, and
Z. Tian, ``Revisiting Weight Regularization for Low-Rank Continual Learning,''
in \emph{International Conference on Learning Representations (ICLR)}, 2026.
[Online]. Available: \url{https://arxiv.org/abs/2602.17559}

\bibitem{shahbazinia2025epismart}
A. Shahbazinia, J. Dan, J. A. Miranda, G. Ansaloni, and D. Atienza,
``Personalization on a Budget: Minimally-Labeled Continual Learning for
Resource-Efficient Seizure Detection,'' 2025.
[Online]. Available: \url{https://arxiv.org/abs/2509.13974}

\bibitem{bassobert2025binaryreplay}
Y. Basso-Bert, A. Molnos, R. Lemaire, W. Guicquero, and A. Dupret,
``Towards Experience Replay for Class-Incremental Learning in Fully-Binary
Networks,'' 2025.
[Online]. Available: \url{https://arxiv.org/abs/2503.07107}

\bibitem{barusco2026vadedge}
M. Barusco, F. Borsatti, D. Petrovic, D. Dalle Pezze, and G. A. Susto,
``Continual Visual Anomaly Detection on the Edge: Benchmark and Efficient
Solutions,'' 2026.
[Online]. Available: \url{https://arxiv.org/abs/2604.06435}

\bibitem{pavan2026odlsurvey}
M. Pavan, L. Pezzarossa, F. Pittorino, M. Roveri, and X. Fafoutis,
``What Changes After Deployment? A Survey on On-Device Learning in TinyML,''
2026. [Online]. Available: \url{https://arxiv.org/abs/2605.31226}

\end{thebibliography}
